%% toc: P3 25 25.4
\section{Saturated elementary extensions of arbitrary structures}
\label{sec:p03-w03-p03-010-saturated-extension}
\index{saturation!existence of saturated extensions}

\paragraph{Purpose.}
Saturated structures are where types become concrete: a type over a small
parameter set that is consistent with the structure is realized by actual
elements.  Vaught's countable construction
(\cref{thm:p03-vaught-saturated-existence}) produces such a model only for a
small countable complete theory and only over finite parameter sets.  This
section proves, for an arbitrary structure \(\mathcal M\) in an arbitrary
first-order language and an arbitrary infinite cardinal \(\kappa\), that
\(\mathcal M\) has a \(\kappa\)-saturated elementary extension.  The proof
supplies the argument behind \cref{thm:saturated-extension} and the
existence clause used in the example after
\cref{thm:saturation-homogeneity}.

\paragraph{Inputs and conclusions.}
The input is a structure \(\mathcal M\) for a first-order language
\(\Lang\) of any cardinality, and an infinite cardinal \(\kappa\).  No
theory is fixed in advance and \(\mathcal M\) need not be infinite.  The
conclusion is an elementary extension \(\mathcal N\succcurlyeq\mathcal M\),
with \(\mathcal M\) an actual elementary substructure, in which every complete
type in finitely many variables over every parameter set of cardinality at
most \(\kappa\) that is consistent with \(\mathcal N\) is realized.  Thus
\(\mathcal N\) is \(\kappa^+\)-saturated and in particular
\(\kappa\)-saturated.  No upper bound on the cardinality of \(\mathcal N\) is
asserted; \cref{thm:p03-w03-p03-010-size} gives the lower bound that every
such infinite structure must meet.

\paragraph{Definitions and assumptions.}
Types over a parameter set \(A\), consistency with a structure, completeness
and realization are those of \cref{sec:types-ultraproducts-01}, using the
parameter expansion \(\mathcal M_A\) of \cref{def:parameter-expansion}.
Elementary diagrams are those of \cref{def:diagrams}.  The ambient set
theory includes the Axiom of Choice; it is used in compactness for a
possibly uncountable language, in cardinal arithmetic, and in choosing the
stages of the construction.  Cardinals are initial ordinals.

\begin{hbconvention}{def:p03-w03-p03-010-saturation-scope}{The three scopes in \(\kappa\)-saturation}
The notion of \(\kappa\)-saturation used throughout is the one defined in
\cref{sec:types-ultraproducts-01}, just before \cref{thm:saturated-extension}:
for an infinite cardinal \(\kappa\), a structure \(\mathcal N\) is
\emph{\(\kappa\)-saturated} when, for every \(A\subseteq\dom\mathcal N\)
with \(\abs A<\kappa\), every complete type over \(A\) in finitely many
variables that is consistent with \(\mathcal N\) is realized in
\(\mathcal N\).  That definition fixes three scopes separately.  The \emph{tuple
scope} is finite: a type is in variables \(x_1,\ldots,x_n\) for some
\(n<\omega\).  The \emph{parameter scope} is a set \(A\) of elements of
\(\mathcal N\) itself, named by the constants of \(\Lang(A)\).  The
\emph{cardinal scope} is the strict bound \(\abs A<\kappa\).  For
\(\kappa=\aleph_0\) the parameter sets are finite; for \(\kappa^+\), the
least cardinal greater than \(\kappa\), they are the sets of cardinality at
most \(\kappa\).  If \(\lambda\leq\mu\), every \(\mu\)-saturated structure is
\(\lambda\)-saturated, because each set of cardinality below \(\lambda\) has
cardinality below \(\mu\).
\index{saturation!tuple, parameter and cardinal scope}
\end{hbconvention}

\subsection{Types under elementary extension}
\label{sec:p03-w03-p03-010-type-facts}

\begin{hblemma}{thm:p03-w03-p03-010-type-facts}{proved=annex:p03-w03-p03-010-type-facts}{Consistency, completion and transfer of types}
Assume the Axiom of Choice.  Let \(\mathcal M\) be an \(\Lang\)-structure,
\(A\subseteq\dom\mathcal M\) and \(p(\bar x)\) a set of
\(\Lang(A)\)-formulas with free variables among
\(\bar x=(x_1,\ldots,x_n)\).
\begin{enumerate}
  \item The set \(p\) is a type over \(A\) consistent with \(\mathcal M\)
        if and only if every finite subset of \(p\) is realized in
        \(\mathcal M\).
  \item Every type over \(A\) consistent with \(\mathcal M\) is contained
        in a complete type over \(A\) consistent with \(\mathcal M\).
  \item If \(\mathcal M\preccurlyeq\mathcal M'\), then
        \(\operatorname{Th}_{\Lang(A)}(\mathcal M_A)=\operatorname{Th}_{\Lang(A)}(\mathcal M'_A)\).
        Hence \(p\) is a (complete) type over \(A\) consistent with
        \(\mathcal M\) exactly when it is one consistent with
        \(\mathcal M'\), and a tuple of \(\mathcal M\) realizes \(p\) in
        \(\mathcal M\) exactly when it realizes \(p\) in \(\mathcal M'\).
\end{enumerate}
\end{hblemma}

\subsection{The one-step extension}
\label{sec:p03-w03-p03-010-one-step}

\begin{hblemma}{thm:p03-w03-p03-010-copy}{proved=annex:p03-w03-p03-010-copy}{Replacing an elementary embedding by an elementary extension}
Let \(h:\mathcal M\to\mathcal N^*\) be an elementary embedding of
\(\Lang\)-structures, and put
\(s_{\mathcal M}=\{u\in\bigcup\bigcup\dom\mathcal M:u\notin u\}\).
There is an \(\Lang\)-structure \(\mathcal N\) with
\(\mathcal M\preccurlyeq\mathcal N\) and an isomorphism
\(F:\mathcal N^*\to\mathcal N\) with \(F(h(a))=a\) for every
\(a\in\dom\mathcal M\).  Its domain is
\[
  \dom\mathcal N=\dom\mathcal M\cup
  \{(y,s_{\mathcal M}):y\in\dom\mathcal N^*\setminus h[\dom\mathcal M]\},
\]
and the two displayed parts are disjoint.
\end{hblemma}

\begin{hblemma}{thm:p03-w03-p03-010-one-step}{proved=annex:p03-w03-p03-010-one-step}{Realizing every type over every subset in one elementary extension}
Assume the Axiom of Choice.  Let \(\mathcal M\) be an \(\Lang\)-structure.
There is an elementary extension \(\mathcal M'\succcurlyeq\mathcal M\)
such that, for every \(A\subseteq\dom\mathcal M\) and every complete type
\(p\) over \(A\) in finitely many variables consistent with \(\mathcal M\),
some tuple of \(\mathcal M'\) realizes \(p\) in \(\mathcal M'\).
Moreover \(\mathcal M'\) can be chosen with
\[
 \dom\mathcal M'\subseteq\dom\mathcal M\cup
 \bigl(\mu_{\mathcal M}\times\{s_{\mathcal M}\}\bigr),
\]
where \(s_{\mathcal M}\) is the set of \cref{thm:p03-w03-p03-010-copy} and
\(\mu_{\mathcal M}\) is the cardinal
\[
 \mu_{\mathcal M}=\max(\abs\Lang,\abs{\dom\mathcal M},\abs{P_{\mathcal M}},\aleph_0),
\]
with \(P_{\mathcal M}\) the set of all pairs \((A,p)\) just described.
The class of all such \(\mathcal M'\) is therefore a nonempty set
determined by \(\mathcal M\).
\end{hblemma}

The one-step lemma realizes types over \emph{all} subsets of
\(\mathcal M\), small ones included.  It cannot by itself give saturation,
because the new elements of \(\mathcal M'\) create new parameter sets whose
types the step did not consider.  The chain below repeats the step, and its
length is chosen so that every small parameter set of the final union has
already appeared at some earlier stage.

\subsection{Regular-length bookkeeping}
\label{sec:p03-w03-p03-010-bookkeeping}

\begin{hblemma}{thm:p03-w03-p03-010-bookkeeping}{proved=annex:p03-w03-p03-010-bookkeeping}{A small set lies in an earlier stage of a chain of length \(\kappa^+\)}
Assume the Axiom of Choice.  Let \(\kappa\) be an infinite cardinal and
\(\langle X_\alpha:\alpha<\kappa^+\rangle\) a sequence of sets with
\(X_\alpha\subseteq X_\beta\) whenever \(\alpha\leq\beta<\kappa^+\).
Then \(\kappa^+\) is a limit ordinal, every \(\alpha<\kappa^+\) has
\(\abs\alpha\leq\kappa\), and every
\(Y\subseteq\bigcup_{\alpha<\kappa^+}X_\alpha\) with \(\abs Y\leq\kappa\)
satisfies \(Y\subseteq X_\gamma\) for some \(\gamma<\kappa^+\).
\end{hblemma}

The lemma is the regularity of \(\kappa^+\) in the form the construction
uses; the general cofinality facts are those of
\cref{thm:p06-cofinality-facts}.  The annex proof uses only
\(\kappa\cdot\kappa=\kappa\) from \cref{thm:p06-cardinal-arithmetic-laws}.

\begin{hbmethod}{def:p03-w03-p03-010-construction}{Building a \(\kappa\)-saturated elementary extension}
Given \(\mathcal M\) and an infinite cardinal \(\kappa\):
\begin{enumerate}
  \item Put \(\mathcal N_0=\mathcal M\).
  \item At a successor stage, let \(\mathcal N_{\alpha+1}\) be an
        extension of \(\mathcal N_\alpha\) given by
        \cref{thm:p03-w03-p03-010-one-step}.  Side condition: the
        extension is chosen from the set of all such extensions with domain
        inside the set that lemma displays, so that one well-ordering made
        in advance can select every stage.
  \item At a limit stage \(\delta\leq\kappa^+\), let
        \(\mathcal N_\delta=\bigcup_{\alpha<\delta}\mathcal N_\alpha\).
        Side condition: \(\langle\mathcal N_\alpha:\alpha<\delta\rangle\)
        must already be an elementary chain, so that
        \cref{thm:elementary-chain} applies.
  \item Output \(\mathcal N=\mathcal N_{\kappa^+}\).  Side condition: the
        chain length is \(\kappa^+\), not \(\omega\) or \(\kappa\); this is
        what \cref{thm:p03-w03-p03-010-bookkeeping} needs.
\end{enumerate}
\end{hbmethod}

\subsection{The existence theorem}
\label{sec:p03-w03-p03-010-existence}

\begin{hbtheorem}{thm:p03-w03-p03-010-saturated-extension}{proved=annex:p03-w03-p03-010-saturated-extension}{Every structure has a \(\kappa\)-saturated elementary extension}
Assume the Axiom of Choice.  Let \(\mathcal M\) be an \(\Lang\)-structure
and \(\kappa\) an infinite cardinal.  The construction of
\cref{def:p03-w03-p03-010-construction} can be carried out, and its output
\(\mathcal N\) satisfies \(\mathcal M\preccurlyeq\mathcal N\).  Every
complete type in finitely many variables over a set
\(A\subseteq\dom\mathcal N\) with \(\abs A\leq\kappa\), consistent with
\(\mathcal N\), is realized in \(\mathcal N\).  Hence \(\mathcal N\) is
\(\kappa^+\)-saturated and \(\kappa\)-saturated.  In particular
\cref{thm:saturated-extension} holds.
\end{hbtheorem}

\begin{hbcorollary}{thm:p03-w03-p03-010-size}{proved=annex:p03-w03-p03-010-size}{Size of saturated structures and the finite case}
Assume the Axiom of Choice, and let \(\kappa\) be an infinite cardinal.
\begin{enumerate}
  \item Every infinite \(\kappa\)-saturated structure has cardinality at
        least \(\kappa\).
  \item Every finite structure is \(\kappa\)-saturated, and every
        elementary extension of a finite structure is that structure
        itself.
\end{enumerate}
\end{hbcorollary}

\begin{hbcorollary}{thm:p03-w03-p03-010-long-tuples}{proved=annex:p03-w03-p03-010-long-tuples}{Realizing types in at most \(\kappa\) variables}
Assume Choice. Let \(\mathcal N\) be \(\kappa\)-saturated in the
finite-tuple sense, let \(A\subseteq\dom\mathcal N\) have size less
than \(\kappa\), and let \(\mu\leq\kappa\) be an ordinal. A complete
type \(p((x_\alpha)_{\alpha<\mu})\) over \(A\), consistent with
\(\mathcal N\), is realized in \(\mathcal N\). Here formulas still
have only finitely many free variables: completeness means deciding every
such formula whose variables belong to the displayed family, and
consistency is joint consistency with
\(\operatorname{Th}_{\Lang(A)}(\mathcal N_A)\).

Consequently allowing fewer than \(\kappa\) variables, with this
finite-formula convention and the same parameter bound, gives an
equivalent definition of \(\kappa\)-saturation. It does not impose a
stronger saturation requirement. Arbitrarily long tuples are different:
an infinite pure equality structure of cardinality \(\kappa\) is
\(\kappa\)-saturated but cannot realize the consistent type of
\(\kappa^+\) pairwise distinct variables.
\end{hbcorollary}


\subsection{Worked example}
\label{sec:p03-w03-p03-010-example}

\begin{hbtheorem}{thm:p03-w03-p03-010-omega-example}{proved=annex:p03-w03-p03-010-omega-example}{An \(\aleph_1\)-saturated elementary extension of \((\N,<)\)}
Assume the Axiom of Choice.  Apply
\cref{thm:p03-w03-p03-010-saturated-extension} to
\(\mathcal M=(\N,<)\) with \(\kappa=\aleph_0\), obtaining an
\(\aleph_1\)-saturated \(\mathcal N\succcurlyeq(\N,<)\).  Call an element of
\(\mathcal N\) \emph{nonstandard} when it is not in \(\N\).
\begin{enumerate}
  \item Every nonstandard element is greater than every natural number,
        and has an immediate predecessor in \(\mathcal N\), which is again
        nonstandard.
  \item Nonstandard elements exist.
  \item For every countable set \(E\) of nonstandard elements there is a
        nonstandard element below every member of \(E\).  Hence no
        countable set of nonstandard elements is coinitial among them.
  \item The domain of \(\mathcal N\) is uncountable.
\end{enumerate}
\end{hbtheorem}

In clause~3 the parameter set is \(\N\cup E\), which is countable and
therefore of cardinality below \(\aleph_1\).  The type used is
\(\{c_n<x:n\in\N\}\cup\{x<c_e:e\in E\}\).  The structure \((\N,<)\)
itself omits even the part \(\{c_n<x:n\in\N\}\) over the countable set
\(\N\), so it is not \(\aleph_1\)-saturated.  Clause~4 is clause~1 of \cref{thm:p03-w03-p03-010-size} for
\(\kappa=\aleph_1\).

\subsection{Limitations}
\label{sec:p03-w03-p03-010-limits}

\begin{hbwarning}{def:p03-w03-p03-010-limits}{What the existence theorem does not give}
\begin{enumerate}
  \item \emph{No size control.}  The theorem bounds the size of
        \(\mathcal N\) from neither side beyond
        \cref{thm:p03-w03-p03-010-size}.  It does not produce a
        \emph{saturated} structure, that is, one that is
        \(\abs{\dom\mathcal N}\)-saturated.  For
        \(\kappa>\aleph_0\) no countable infinite structure is
        \(\kappa\)-saturated, by that corollary.
  \item \emph{The chain length matters to the argument.}  For a chain of
        length \(\omega\) the bookkeeping step can fail.  If the
        domains \(X_n\) increase strictly, choosing
        \(y_n\in X_{n+1}\setminus X_n\) gives a countable set
        \(\{y_n:n<\omega\}\) lying in no \(X_n\).  In the construction
        with infinite \(\mathcal M\), each step does add elements, since
        the type \(\{x\ne c_a:a\in\dom\mathcal N_n\}\) over the whole
        stage completes to a type that the next stage realizes.  So the
        proof given here does not show that a union of length \(\omega\)
        is \(\aleph_1\)-saturated; this records a failure of the argument,
        not a theorem that such a union is never \(\aleph_1\)-saturated.
  \item \emph{Tuple length.} The construction directly realizes types
        in finitely many variables. By
        \cref{thm:p03-w03-p03-010-long-tuples}, this already realizes
        types in at most \(\kappa\) variables over fewer than
        \(\kappa\) parameters. Formulas still involve finitely many
        variables each. There is no corresponding guarantee for arbitrary
        larger variable families, as the pure equality example shows.
  \item \emph{Choice.}  Compactness for the language \(\Lang(M)\), which
        may be uncountable, the cardinal arithmetic, and the selection of
        stages all use the Axiom of Choice.
\end{enumerate}
\end{hbwarning}

\paragraph{Relationships.}
\Cref{thm:p03-w03-p03-010-saturated-extension} with
\cref{thm:saturation-homogeneity} shows that every structure has a
\(\kappa\)-homogeneous elementary extension for every infinite \(\kappa\).
Vaught's countable theorem
(\cref{thm:p03-vaught-saturated-existence,thm:p03-vaught-smallness-equivalences})
is sharper in one respect: it yields a \emph{countable} model saturated
over finite parameter sets, exactly when the countable complete theory is
small.  The present theorem needs no hypothesis on the theory or the
language, and pays for this with the loss of any cardinality bound.  For
\(\aleph_0\)-saturation alone, where the parameter sets are finite, the
regular length is not needed: each element of a finite subset of the union
of an increasing chain of length \(\omega\) lies in some stage, and the
largest of these finitely many stages contains the whole subset.  The same
steps along a chain of length \(\omega\) therefore give an
\(\aleph_0\)-saturated elementary extension.  The length \(\kappa^+\) and
\cref{thm:p03-w03-p03-010-bookkeeping} are what infinite parameter sets
require, and clause~2 of \cref{def:p03-w03-p03-010-limits} records why the
argument does not reach them with the shorter chain.

\paragraph{Sources.}
Morley and Vaught characterize their homogeneous-universal systems of
power \(\kappa\) by realization, in an expansion naming a set of fewer
than \(\kappa\) elements, of every finitely satisfiable set of formulas in
one free variable \cite[p.~49, Theorem~3.4]{MorleyVaught1962}; their
\(S_\kappa(X)\) is the family of subsets of \(X\) of power less than
\(\kappa\) \cite[p.~39]{MorleyVaught1962}, and their systems are
relational.  Their existence
theorem for such a system of power exactly \(\kappa\) carries additional
cardinal and type-counting hypotheses
\cite[p.~49, Theorem~3.5]{MorleyVaught1962}, with \(\kappa^*\) defined
on p.~38 there; the present theorem drops those hypotheses together with
the size conclusion.  The one-step compactness lemma and the passage to an
elementary chain indexed by a regular cardinal follow van den Berg's
course notes for the 2015--2016 Mastermath model theory course, chapter~10,
a lecture handout rather than a published text
\cite[ch.~10, pp.~2--3, Theorem~10.4 and Lemma~10.5]{vandenBerg2016ModelTheoryCh10}.
There saturation is defined with one-variable types
\cite[ch.~10, p.~1, Definition~10.1]{vandenBerg2016ModelTheoryCh10}, the
chain argument is written out for chains of length \(\omega\), and the
regularity of \(\kappa^+\) is stated without proof.  The version here realizes complete types in finitely many
variables at each step, makes the inclusion \(\mathcal M\subseteq\mathcal N\)
actual by \cref{thm:p03-w03-p03-010-copy}, proves the regularity step in
\cref{thm:p03-w03-p03-010-bookkeeping} and justifies the choice of the
stages inside a set.
